% Draft of Section 2 (Data Exploration). Requires \usepackage{graphicx,booktabs}.
% All numbers come from data_exploration.py (tables/ and figures/).

\section{Data Exploration}

\paragraph{Features and class balance.}
The data has 870 rows from 290 respondents; each respondent answered the same
12 questions once per sport, so Baseball, Basketball and Soccer have exactly 290
rows each and no rebalancing is needed. Features: one numeric rating
(\emph{ease}, 0--10), three 5-level Likert items (\emph{strategy},
\emph{injury}, \emph{individual}), three multi-select questions (\emph{where},
\emph{season}, \emph{music}; 21 options in total), four free-text answers
(\emph{food}, \emph{colour}, \emph{song}, \emph{feel}) and a free-text price
(\emph{pay}). The respondent \texttt{id} is not a feature: it says nothing about
the sport and test respondents never appear in training.

\paragraph{Test set.}
Before any exploration we shuffled the respondent ids (seed 311) and split them
70/15/15 into train (203 respondents, 609 rows), validation (44, 132) and test
(43, 129), each balanced across classes. We split by respondent because a
person's three rows are dependent: the id alone explains 39\% of the variance of
\emph{ease}. All statistics, vocabularies, imputation and scaling values below
come from the training split only; validation and test rows were never
inspected.

\paragraph{Missing values and inconsistencies.}
Besides empty strings, missing answers appear as \emph{none}, \emph{n/a},
\emph{idk}, \emph{no idea}, \emph{nothing in particular}, \emph{I don't watch
this sport} and the Excel error \texttt{\#NAME?}; we map all of them to one
marker (Table~\ref{tab:missing}). Seven training rows are completely empty and
are dropped. \emph{song} is missing far more often for Baseball and Basketball
than for Soccer (almost everyone recalls \emph{Waka Waka}), so missingness is
predictive and we add a binary \emph{song missing} flag. \emph{No particular
type} (\emph{music}) is a real option and \emph{none} in \emph{pay} means \$0, so
neither counts as missing. Two \emph{ease} values are non-integers (3.5, 7.5) and
are kept. Remaining gaps are imputed with the training median (numeric) or left
as all-zero indicators (multi-select, text).

\begin{table}[htbp]
\centering\small
\begin{tabular}{lcccccc}
\toprule
Missing (\%) & ease, Likert, multi-select & food & colour & song & feel & pay \\
\midrule
Overall    & 1.1--1.3 & 4.6 & 2.0 & 18.2 & 4.4 & 3.0 \\
Baseball   & 1.5--2.0 & 3.9 & 2.5 & 26.1 & 6.4 & -- \\
Basketball & 1.5--2.0 & 8.4 & 2.0 & 22.2 & 3.9 & -- \\
Soccer     & 0.5      & 1.5 & 1.5 & 6.4  & 3.0 & -- \\
\bottomrule
\end{tabular}
\caption{Missing values in the training split after normalizing missing markers.}
\label{tab:missing}
\end{table}

\paragraph{Parsing \emph{pay}; outliers.}
Prices come in many formats (\emph{\$1,000}, \emph{200CAD}, \emph{1k},
\emph{Thirty dollars}, \emph{\$150--\$250}, \emph{At most 500 CAD.}). We remove
thousands separators, expand \emph{k} and spelled-out numbers, take the midpoint
of a range and otherwise the first number; clear refusals (\emph{not willing},
\emph{nothing}) become 0. Answers on a ``1 to 10 scale'' or without a number
(\emph{very willing}, \emph{hundreds of dollars}) are unparseable (1.8\%) and
treated as missing. Values are heavily right-skewed (median \$100, mean \$262,
max \$10{,}000; 12 answers exceed the 1.5$\times$IQR fence on the log scale).
High prices are plausible for finals, so we cap at the training 99th percentile
(\$3{,}000) and use $\log(1+\textit{pay})$.

\paragraph{Numeric and Likert features.}
Figure~\ref{fig:ordinal} shows monotone shifts in class means: Baseball is rated
harder to start (mean \emph{ease} 4.06 vs.\ 6.06/6.22), less risky and more
individual. The gap between the most distant class means divided by the pooled
within-class SD is 1.02 (\emph{ease}), 0.87 (\emph{injury}), 0.82
(\emph{individual}), 0.61 (\emph{pay}) and 0.48 (\emph{strategy}): useful, but
mostly for separating Baseball from the rest. Spearman correlations between
these features are at most 0.30, so none is redundant. Because the effect is
monotone, Likert answers are encoded as ordinal values 1--5 (one column instead
of five one-hot columns).

\begin{figure}[htbp]
\centering
\includegraphics[width=0.92\linewidth]{figures/05_numeric_likert_by_class.png}
\caption{\emph{ease} per class (left) and Likert answers per class (1 = strongly
disagree, 5 = strongly agree). Training split.}
\label{fig:ordinal}
\end{figure}

\paragraph{Multi-select features.}
Each option becomes a binary column. Several options separate classes sharply:
\emph{Winter} (Basketball 59\% vs.\ $\sim$5\%), \emph{Indoor facility}
(Basketball 96\% vs.\ Baseball 26\%), \emph{Latin} (Soccer 44\% vs.\ $\le$4\%),
\emph{Hip-hop or rap} (Basketball 84\%) and \emph{Country} (Baseball 32\%).
\emph{Summer} is ticked by 94\% of every class and is dropped as uninformative.

\paragraph{Free-text features.}
Answers are short (median 1--4 words). We lower-case them, remove the
placeholders \emph{[this sport]}/\emph{[a competition]} and stop words, and
build a binary bag-of-words per question on the training split, keeping words
in at least 3 answers (95/25/94/120 words for food/colour/song/feel); unseen
words are ignored at prediction time. A few words are nearly
class-pure (Figure~\ref{fig:text}): \emph{orange} (111/116 Basketball),
\emph{green} (99/110 Soccer), \emph{waka} (101/101 Soccer), \emph{ball game}
(Baseball), \emph{drake} (Basketball), \emph{relaxed}/\emph{bored} (Baseball);
most others are weak. Since counts are almost always 0/1, TF-IDF would mainly
re-weight rare words; we compare both in Methodology.

\begin{figure}[htbp]
\centering
\includegraphics[width=0.92\linewidth]{figures/09_text_words_by_class.png}
\caption{Ten most frequent words per text question: counts per class; shade =
share of the word's occurrences in that class.}
\label{fig:text}
\end{figure}

\paragraph{Scaling, dimensionality and predictive features.}
The encoding gives about 361 columns for 602 training rows, so we favour simple
models with strong regularization. Continuous features are standardized with
training mean and SD (reused on validation/test to avoid leakage); this matters
for k-NN and softmax regression with L2, not for trees or Naive Bayes. Training
a softmax regression on one feature group at a time (5-fold CV grouped by
respondent, training split) gives accuracies of 0.75 (\emph{colour}), 0.71
(\emph{song}), 0.69 (\emph{music}), 0.60 (\emph{where}), 0.56 (\emph{food}),
0.53 (\emph{season}) and 0.43--0.49 for the rest (chance 0.33). Many sparse
indicators that each add a little evidence suit Naive Bayes and regularized
softmax regression; sharp thresholds (\emph{Winter}, \emph{Indoor},
\emph{ease}) suit decision trees and random forests.
